\chapter*{第 8 章 黏性流体动力学基础}\addcontentsline{toc}{chapter}{第 8 章 黏性流体动力学基础}\markboth{第 8 章 黏性流体动力学基础}{}

\begin{tishi}
\textbf{章号与题源说明：}本册按\textbf{赵琴教材第 1--8 章}归章，本章为\textbf{赵琴第 8 章 黏性流体动力学基础}。
孔珑原书的对应关系是：\textbf{孔珑第 9 章 黏性流体绕过物体的流动 $\to$ 赵琴第 8 章}，
对应主题为边界层的基本概念与厚度、普朗特边界层方程、边界层动量积分方程、
平板边界层计算（层流／湍流／混合）、边界层分离与形状阻力、绕流阻力与减阻。
本章收录孔珑第 9 章课后习题中的
\textbf{9-9、9-10、9-11、9-12、9-14、9-15、9-16、9-17、9-18、9-19、9-20 共 11 题}；
其余题号略去的理由见本章末说明。

\textbf{本册星级规则：}5 星＝该考点在历年真题中作为计算题出现 3 次以上，或考情分析列为计算题重点；
4 星＝真题中出现 2--3 次；3 星＝真题中出现 1 次，或考情分析只要求"掌握基本概念、学会套用公式"；
2 星＝真题未出现但属考纲核心内容；1 星＝仅作扩展了解。

\textbf{做题要求：}每题必须写出"已知—公式—代入—单位—结果"五步，单位一律化为国际单位制；
边界层题先算 $\Rey_x$ 或 $\Rey_l$ 判断流态，是混合边界层的必须先求出转捩点位置 $x_{\mathrm{cr}}$。
\end{tishi}

\begin{ti}
\sloppy
\emergencystretch=2em

\item \textbf{孔珑 9-9}\quad 温度为 $25\ ^{\circ}\mathrm{C}$ 的空气，以 $30\ \mathrm{m/s}$ 的速度纵向绕流一块极薄的平板，压强为大气压，计算离平板前缘 $200\ \mathrm{mm}$ 处边界层的厚度为多少？
\tuyi{极薄平板顺来流放置，来流速度 $v_\infty=30\ \mathrm{m/s}$；平板前缘为边界层的起点，在 $x=200\ \mathrm{mm}$ 处量取边界层厚度 $\delta$。}
\quad\textbf{重要度}\xstarfull{4}\quad\yiju{E4 2020 填空 6（油绕 $L=2\ \mathrm{m}$ 平板求末端边界层厚度，3 分）、E4 2022 填空 6（水绕 $L=2\ \mathrm{m}$ 平板判断末端流态，3 分）、E4 2021 计算 5（平板尾端边界层厚度与两面摩擦阻力，10 分）；E1 slide33 第八章"会涉及选择题和填空题考点"、E2 slide5 第六至八章"掌握基本概念、学会套用计算公式"；E6 教材 8.2、8.4（$\delta=5.84x/\sqrt{\Rey_x}$ 属布拉修斯解口径）}\quad\nandu{易}

\item \textbf{孔珑 9-10}\quad 温度为 $20\ ^{\circ}\mathrm{C}$、密度为 $925\ \mathrm{kg/m^3}$、运动黏度为 $7.9\times10^{-5}\ \mathrm{m^2/s}$ 的油流，以 $60\ \mathrm{cm/s}$ 的速度纵向绕流一长 $50\ \mathrm{cm}$、宽 $15\ \mathrm{cm}$ 的薄平板。试求总摩擦阻力和边界层厚度。
\tuyi{薄平板顺油流放置，板长 $L=50\ \mathrm{cm}$、板宽 $b=15\ \mathrm{cm}$，来流速度 $v_\infty=60\ \mathrm{cm/s}$；平板两面都与油接触，故总摩擦阻力为单面的 2 倍，边界层厚度按板末端 $x=L$ 计算。}
\quad\textbf{重要度}\xstarfull{5}\quad\yiju{E4 2015 计算 9（$2\ \mathrm{m}\times8\ \mathrm{m}$ 平板摩擦阻力之比 $F_1/F_2$，12 分）、E4 2016 计算 10（船体按相当平板求摩擦阻力与功率，10 分）、E4 2018 计算 9（平底船底板边界层阻力功率，10 分）、E4 2021 计算 5（平板两面摩擦阻力，10 分），平板摩擦阻力已在计算题中出现 4 次；E6 教材 8.4、8.6}\quad\nandu{中}

\item \textbf{孔珑 9-11}\quad 平板层流边界层内速度分布规律为 $v_x/v_\infty=2\,(y/\delta)-(y/\delta)^{2}$，试求：
\begin{xiaowen}
\item 边界层厚度 $\delta$ 与雷诺数 $\Rey$ 的关系式；
\item 摩擦阻力系数 $C_{\mathrm{Df}}$ 与雷诺数 $\Rey$ 的关系式。
\end{xiaowen}
\quad\textbf{重要度}\xstarfull{3}\quad\yiju{E4 2021 计算 5 卷面直接给出与本题同源的层流边界层公式 $\delta/x=5.48/\sqrt{\Rey_x}$（即本题抛物线速度剖面经动量积分方程得到的结果）与 $C_f=1.328/\sqrt{\Rey_l}$；E6 教材 8.5（冯·卡门动量积分方程）；E2 slide5 第六至八章"学会套用计算公式即可"}\quad\nandu{难}

\item \textbf{孔珑 9-12}\quad 若平板层流边界层内的速度分布为正弦曲线 $v_x=v_\infty\sin\dfrac{\pi y}{2\delta}$，试求边界层厚度 $\delta$、摩擦阻力系数 $C_{\mathrm{Df}}$ 与雷诺数 $\Rey$ 之间的关系式。
\quad\textbf{重要度}\xstarfull{3}\quad\yiju{E4 2021 计算 5 卷面给出的层流边界层公式族（$\delta/x=5.48/\sqrt{\Rey_x}$、$C_f=1.328/\sqrt{\Rey_l}$）与本题同类，都是用假定速度剖面代入动量积分方程的结果；E6 教材 8.5；E2 slide5 第六至八章"学会套用计算公式即可"}\quad\nandu{难}

\item \textbf{孔珑 9-14}\quad 在相同雷诺数 $\Rey_l$ 的情况下，试求 $20\ ^{\circ}\mathrm{C}$ 的水和 $30\ ^{\circ}\mathrm{C}$ 的空气各平行流过长度为 $l$ 的平板时产生的摩擦阻力之比。
\tuyi{同一块长 $l$、宽 $b$ 的平板，分别放在 $20\ ^{\circ}\mathrm{C}$ 的水流和 $30\ ^{\circ}\mathrm{C}$ 的空气流中，两种情况下板末端雷诺数 $\Rey_l$ 相同；因此两种流体的来流速度并不相同（由 $\Rey_l$ 反推）。}
\quad\textbf{重要度}\xstarfull{4}\quad\yiju{E4 2015 计算 9（同一平板两种放置方式的摩擦阻力之比 $F_1/F_2$，12 分，同一"阻力之比"题型）；E4 2016 计算 10、E4 2018 计算 9、E4 2021 计算 5（平板或相当平板的摩擦阻力）；E6 教材 8.6}\quad\nandu{中}

\item \textbf{孔珑 9-15}\quad 一块长 $6\ \mathrm{m}$、宽 $2\ \mathrm{m}$ 的平板平行静止地安放在速度为 $60\ \mathrm{m/s}$ 的 $40\ ^{\circ}\mathrm{C}$ 空气流中，在平板边界层内从层流转变为紊流的临界雷诺数 $\Rey_{x\mathrm{cr}}=10^{6}$。试计算平板的摩擦阻力。
\tuyi{平板顺气流放置，板长 $L=6\ \mathrm{m}$、板宽 $b=2\ \mathrm{m}$；气流先在前缘一段内形成层流边界层，到 $x=x_{\mathrm{cr}}$ 处转捩为湍流边界层，因此是"层流段＋湍流段"的混合边界层，需分别处理再相减。}
\quad\textbf{重要度}\xstarfull{5}\quad\yiju{E4 2016 计算 10（船体按相当平板求混合边界层摩擦阻力与功率，10 分）、E4 2018 计算 9（船底平板混合边界层阻力功率，10 分）、E4 2021 计算 5（平板混合边界层厚度与两面摩擦阻力，10 分），混合边界层阻力已在计算题中出现 3 次；E6 教材 8.6}\quad\nandu{难}

\item \textbf{孔珑 9-16}\quad 直径为 $500\ \mathrm{mm}$ 的管道，通过 $30\ ^{\circ}\mathrm{C}$ 的空气，在垂直于管道的轴线方向插入直径为 $10\ \mathrm{mm}$ 的卡门涡街流量计，测得漩涡的脱落频率为 $105\ (1/\mathrm{s})$，求管道中的流量。
\tuyi{在直径 $D=500\ \mathrm{mm}$ 的管道中插入直径 $d=10\ \mathrm{mm}$ 的圆柱形涡街发生体，气流绕过该柱体时在柱后交替脱落漩涡，脱落频率 $f=105\ \mathrm{s^{-1}}$ 与柱侧流速成正比。}
\quad\textbf{重要度}\xstarfull{1}\quad\yiju{E4 2014--2022 历年真题与 E5 题库中均未见卡门涡街、斯特劳哈尔数题；E6 教材 8.7 只讲边界层分离与尾涡的形成，未引入斯特劳哈尔数 $Sr$；E1 slide33／E2 slide5 第八章只考选择填空——本题按本册星级表"仅作扩展了解"定为一星}\quad\nandu{中}

\item \textbf{孔珑 9-17}\quad 已知汽车的行驶速度为 $60\ \mathrm{km/h}$，垂直于运动方向的投影面积为 $2\ \mathrm{m^2}$，阻力系数为 $0.3$，静止空气的温度为 $0\ ^{\circ}\mathrm{C}$，试求汽车克服空气阻力所作的功率。
\tuyi{汽车以 $v=60\ \mathrm{km/h}$ 匀速行驶，迎面投影面积（迎风面积）$A=2\ \mathrm{m^2}$；车身受到的空气阻力由钝体绕流的压差阻力（形状阻力）决定，用阻力系数 $C_D=0.3$ 计算。}
\quad\textbf{重要度}\xstarfull{4}\quad\yiju{E4 2017 计算 2（汽车以 $80\ \mathrm{km/h}$ 行驶、迎风面积 $2\ \mathrm{m^2}$、$C_D=0.4$，求克服空气阻力所消耗的功率，8 分，与本题同型）；E4 2019 计算 7（用阻力系数法求绕流阻力）；E6 教材 8.7}\quad\nandu{易}

\item \textbf{孔珑 9-18}\quad 在风洞中，温度为 $20\ ^{\circ}\mathrm{C}$ 的空气以 $10\ \mathrm{m/s}$ 的风速垂直吹向直径为 $50\ \mathrm{cm}$ 的圆盘，试求作用在圆盘上的力。
\tuyi{圆盘垂直于气流方向放置（正面迎风），直径 $d=50\ \mathrm{cm}$；气流绕圆盘流动时在圆盘后侧形成分离涡区，圆盘受到的力主要是压差阻力，迎面投影面积为圆面积 $\pi d^{2}/4$。}
\quad\textbf{重要度}\xstarfull{4}\quad\yiju{E4 2017 计算 2（阻力系数法求绕流阻力）、E4 2019 计算 7（斯托克斯区 $C_D=24/\Rey$ 求小雷诺数绕流阻力）、E4 2020 判断 5（无限长圆柱绕流可简化为平面流动）；E6 教材 8.7（圆柱体阻力系数与雷诺数的关系曲线）}\quad\nandu{中}

\item \textbf{孔珑 9-19}\quad 将密度 $\rho=899.4\ \mathrm{kg/m^3}$ 的透平油放在有刻度的玻璃量筒内，让直径 $3\ \mathrm{mm}$、密度 $\rho_s=7791\ \mathrm{kg/m^3}$ 的小钢球在油中自由降落，用秒表测得等速降落的钢球通过量筒上两个刻度之间的时间，算得自由沉降速度为 $11\ \mathrm{cm/s}$，试按教材中式(9-59)求出油的黏度，并验算一下雷诺数是否符合教材中式(9-59)的适用范围。
\begin{xiaowen}
\item 求透平油的动力黏度 $\mu$；
\item 验算雷诺数 $\Rey$ 是否在斯托克斯公式的适用范围内。
\end{xiaowen}
\quad\textbf{重要度}\xstarfull{4}\quad\yiju{E4 2019 计算 7（小钢球在油中自由沉降测油的黏度：油 $\rho=900\ \mathrm{kg/m^3}$、钢球 $d=3\ \mathrm{mm}$、$v=12\ \mathrm{cm/s}$，10 分，与本题同题）；E4 2017 计算 2（阻力系数法求绕流阻力）；E6 教材 8.7}\quad\nandu{中}

\item \textbf{孔珑 9-20}\quad 某采暖沸腾炉的料层温度为 $1000\ ^{\circ}\mathrm{C}$，烟气的运动黏度为 $\nu=1.67\times10^{-6}\ \mathrm{m^2/s}$，料层内燃料的颗粒密度为 $2294\ \mathrm{kg/m^3}$，平均粒径为 $1.68\ \mathrm{mm}$。试问通过料层最小需要多大风速，才能使颗粒处于悬浮状态？
\tuyi{烟气自下而上吹过料层，颗粒受到向上的阻力 $F_D$、向上的浮力 $F_B$ 和向下的重力 $G$；风速小时颗粒静止，风速增大到 $F_D+F_B=G$ 时颗粒开始悬浮，此时的风速即为所求的最小风速。}
\quad\textbf{重要度}\xstarfull{2}\quad\yiju{E4 2014--2022 历年真题与 E5 题库中均未见料层颗粒悬浮（最小流化风速）题；E6 教材 8.7 给出绕流阻力与雷诺数分区（$\Rey<1$ 取 $C_D=24/\Rey$，$1000<\Rey<2\times10^{5}$ 取 $C_D=0.45$）——本题按本册星级表"真题未出现但属考纲核心（绕流阻力）"定为二星}\quad\nandu{难}

\end{ti}

\begin{tishi}
\textbf{本章略去的题号及原因：}
（1）\textbf{孔珑 9-1 至 9-8}（两平行平板反向运动的切应力、倾斜平面上薄液层的均匀层流、
斜楔形滑块的油膜承载与阻力、两同心圆管间的旋转流动、柱塞与油缸间隙的流量与阻力、
环形管道及控制阀缝隙的漏油量等）都是"缝隙流动"的层流精确解，
落在\textbf{赵琴第 5 章黏性流体的层流流动}与\textbf{赵琴第 8 章 8.1 节纳维--斯托克斯方程}的范围，
按本册归章规则应归入第 5 章（层流流动），本章不重复收录；
（2）\textbf{孔珑 9-13} 原书解答即为"略"（无解答内容），无法收录；
（3）\textbf{孔珑 9-21}（圆形、扁形喷管的自由射流中心风速衰减）属自由射流，
不在赵琴第 1--8 章范围内，故不收。
孔珑第 9 章没有气体动力学（激波、马赫数、喷管）题，也没有流体机械、明渠堰流题，
不存在因这些原因略去的题号。
\end{tishi}

\begin{tishi}
\textbf{OCR 校订清单与存疑说明（已对照原书影印页逐条核对）：}
（1）\textbf{孔珑 9-10}：原书在求摩擦阻力处用 $\Rey_l=3.797\times10^{3}$，
而在求边界层厚度处印成 $\Rey=3.797\times10^{4}$，同题自相矛盾；
按 $\Rey_l=v_\infty L/\nu=0.6\times0.5/(7.9\times10^{-5})=3.797\times10^{3}$ 判断，
\textbf{正确值应为 $3.797\times10^{3}$（层流）}——原书给出的 $\delta=47.387\ \mathrm{mm}$
恰好是用 $3.797\times10^{3}$ 算出的，可见厚度算式中的 $10^{4}$ 系原书笔误或 OCR 误读；
（2）\textbf{孔珑 9-15}：原书求平板末端雷诺数的算式中分子印作 $60\times20$（相当于取板长 $20\ \mathrm{m}$），
得 $\Rey_l=7.1\times10^{7}$；而题给板长为 $6\ \mathrm{m}$，按 $60\times6$ 计算应为 $2.13\times10^{7}$。
本册在第 8.6 题中给出两种口径：按原书取 $\Rey_l=7.1\times10^{7}$ 得 $F_{D\mathrm{f}}=53.285\ \mathrm{N}$，
按题给尺寸取 $\Rey_l=2.13\times10^{7}$ 得 $F_{D\mathrm{f}}=61.2\ \mathrm{N}$；
（3）\textbf{孔珑 9-17}：原书功率 OCR 作 $1.7958\times10^{2}\ \mathrm{W}$，
按 $P=F_Dv=107.75\times16.67$ 应为 $1.7958\times10^{3}\ \mathrm{W}$，本册按 $1.80\ \mathrm{kW}$ 订正；
（4）\textbf{孔珑 9-19}：原书多处把密度单位印成 $\mathrm{kg/m^2}$，一律订正为 $\mathrm{kg/m^3}$；
（5）\textbf{孔珑 9-12}：动量损失厚度精确值为 $\delta_2=(4-\pi)\delta/(2\pi)=0.1366\delta$，
原书取 $0.137\delta$，两者相差仅 $0.3\%$，本册沿用原书的 $0.137\delta$；
（6）\textbf{孔珑 9-9、9-10} 用 $\delta=5.84x\Rey_x^{-1/2}$（布拉修斯解口径），
而 9-11、9-12 由动量积分方程得到的是 $5.48x\Rey_x^{-1/2}$ 与 $4.79x\Rey_x^{-1/2}$，
2021 年真题卷面给出的是 $5.48/\sqrt{\Rey_x}$——三者口径不同，做题时\textbf{以卷面给定公式为准}。
\end{tishi}

\begin{tishi}
\textbf{题面 OCR 不清、核对影印页后补齐的说明：}
（1）\textbf{孔珑 9-11}：题面在 OCR 中开头残缺（只残存末尾"与雷诺数 $\Rey$ 的关系式"），
据原书影印页补全为"平板层流边界层内速度分布规律为 $v_x/v_\infty=2\,(y/\delta)-(y/\delta)^{2}$"，
与原书解答中 $\tau_w=2\mu v_\infty/\delta$、$\delta_2=2\delta/15$ 完全吻合；
（2）\textbf{孔珑 9-17}：题面 OCR 缺"阻力系数为 $0.3$"一处，据原书解答
$F_D=0.3\times\frac{1}{2}\times1.293\times(60\times10^{3}/3600)^{2}\times2=107.75\ \mathrm{N}$ 反推补入；
（3）\textbf{孔珑 9-16}：题面给出空气温度为 $30\ ^{\circ}\mathrm{C}$，原书解答中取
$\nu=15.95\times10^{-6}\ \mathrm{m^2/s}$，与 $30\ ^{\circ}\mathrm{C}$ 空气的运动黏度相符，题面无缺项。
\end{tishi}

\begin{tishi}
\textbf{超出 818 考纲的扩展内容提示（题目仍收录，相应公式可不记）：}
（1）\textbf{孔珑 9-16} 的卡门涡街流量计要用斯特劳哈尔数 $Sr=fd/v$；
赵琴教材第 8 章只讲边界层分离与尾涡，未引入 $Sr$，属扩展内容（仅 1 星）；
（2）\textbf{孔珑 9-20} 的料层颗粒悬浮要按雷诺数大小分区选取阻力系数
（$\Rey<1$ 取 $C_D=24/\Rey$，$1000<\Rey<2\times10^{5}$ 取 $C_D=0.45$），
"先假设、算完再回代验算"的试算步骤超出 818 直接考点，
但钝体绕流阻力公式 $F_D=C_D\frac{1}{2}\rho v_\infty^{2}A$ 必须掌握；
（3）\textbf{孔珑 9-15} 中的常数 $A=3300$ 取自教材表 9-1（对应 $\Rey_{x\mathrm{cr}}=10^{6}$），
若考试给出 $\Rey_{x\mathrm{cr}}=5\times10^{5}$，应改查对应的 $A$ 值或直接用卷面给出的阻力系数公式；
（4）\textbf{孔珑 9-14} 的结论"相同 $\Rey_l$ 时摩擦阻力之比等于 $\mu\nu$ 之比"是层流平板公式的推论，
只作理解性掌握，不必背记。
\end{tishi}
